\section{Learning Objectives}

This chapter begins the advanced track. From here on, all math is hand-rolled —
no \texttt{rand}, \texttt{nalgebra}, or \texttt{statrs}. After completing it
you can:

\begin{itemize}
    \item Build a validated price series that rejects non-finite and
      non-positive values at construction time
    \item Compute the first four sample moments (mean, variance, skewness,
      excess kurtosis) from first principles
    \item Apply any aggregation over a sliding window with a generic
      \texttt{rolling} function
    \item Implement a deterministic PRNG (xorshift64\*) and a normal sampler
      (Box-Muller) without external crates
    \item Simulate geometric Brownian motion paths and understand the exact
      solution to the GBM SDE
    \item Detect fat tails via excess kurtosis and explain why Gaussian risk
      models underestimate tail risk
\end{itemize}

\begin{keyinsight}
The Gaussian distribution has excess kurtosis $0$. Real financial returns
have positive excess kurtosis: large losses and gains happen more often than
a normal model predicts. Any risk model that assumes Gaussian returns will
underestimate the probability of extreme events. The 2008 crisis was in
part a fat-tail crisis.
\end{keyinsight}

\section{A Validated Price Series}

\marginnote{%
\textbf{The invariant}\\[0.3em]
Every element of a \texttt{PriceSeries} is strictly positive and finite. The
constructor rejects the rest.
}

The foundation of every quant computation is a clean price series. We wrap
\texttt{Vec<f64>} in a newtype that enforces the invariant ``strictly positive
and finite'' at construction time:

\begin{lstlisting}[language=Rust]
#[derive(Debug, Clone, PartialEq)]
pub struct PriceSeries(Vec<f64>);

impl PriceSeries {
    pub fn new(prices: Vec<f64>) -> Result<Self, CoreError> {
        for &p in &prices {
            if !p.is_finite() || p <= 0.0 {
                return Err(CoreError::InvalidPrice);
            }
        }
        Ok(Self(prices))
    }
    pub fn as_slice(&self) -> &[f64] { &self.0 }
    pub fn len(&self) -> usize { self.0.len() }
}
\end{lstlisting}

Rejecting \texttt{NaN} and \texttt{+inf} at the boundary means every downstream
function can assume finite inputs and skip defensive checks. This is the
``parse, don't validate'' pattern applied to numeric data.

\section{Returns, Again}

Simple and log returns were introduced in \cref{ch:returns}. We re-implement
them here so \texttt{quant-core} is self-contained:

\begin{lstlisting}[language=Rust]
pub fn simple_returns(prices: &[f64]) -> Vec<f64> {
    if prices.len() < 2 { return Vec::new(); }
    let mut out = Vec::with_capacity(prices.len() - 1);
    for i in 1..prices.len() {
        let prev = prices[i - 1];
        out.push(if prev == 0.0 { 0.0 }
                 else { (prices[i] - prev) / prev });
    }
    out
}
\end{lstlisting}

\begin{keyinsight}
The identity $\ln(1 + r^{\text{simple}}) = r^{\text{log}}$ holds exactly. For
small returns the two are nearly equal; for large returns they diverge. This
identity is a cheap property test that catches implementation bugs in either
function.
\end{keyinsight}

\section{Moments}

\marginnote{%
\textbf{Central moment}\\[0.3em]
$m_k = \dfrac{1}{n}\sum_{i=1}^{n}(x_i - \bar x)^k$\\[0.5em]
\textbf{Skewness}\\[0.3em]
$g_1 = \dfrac{m_3}{m_2^{3/2}}$\\[0.3em]
\textbf{Excess kurtosis}\\[0.3em]
$g_2 = \dfrac{m_4}{m_2^2} - 3$
}

The $k$-th central moment is
\[
m_k = \frac{1}{n}\sum_{i=1}^{n}(x_i - \bar x)^k.
\]
Skewness and excess kurtosis are standardised ratios of these moments:

\begin{align*}
g_1 &= \frac{m_3}{m_2^{3/2}} \quad \text{(skewness)} \\
g_2 &= \frac{m_4}{m_2^2} - 3 \quad \text{(excess kurtosis)}
\end{align*}

The Gaussian distribution has $g_1 = 0$ and $g_2 = 0$. Positive skewness means
a longer right tail; negative skewness a longer left tail. Positive excess
kurtosis means heavier tails than Gaussian.

\begin{lstlisting}[language=Rust]
pub fn skewness(data: &[f64]) -> Result<f64, CoreError> {
    if data.len() < 3 { return Err(CoreError::InsufficientData { .. }); }
    let m = mean(data);
    let m2 = central_moment_2(data, m);
    let m3 = central_moment_3(data, m);
    Ok(m3 / m2.powf(1.5))
}

pub fn excess_kurtosis(data: &[f64]) -> Result<f64, CoreError> {
    if data.len() < 4 { return Err(CoreError::InsufficientData { .. }); }
    let m = mean(data);
    let m2 = central_moment_2(data, m);
    let m4 = central_moment_4(data, m);
    Ok(m4 / (m2 * m2) - 3.0)
}
\end{lstlisting}

\begin{keyinsight}
We use the sample variance ($n - 1$ denominator) for $\sigma^2$, consistent
with \func{qf\_common::compute\_stats}, but use population central moments
(denominator $n$) for $m_2, m_3, m_4$ in skewness and kurtosis. This is the
textbook definition; the difference between the two vanishes for large $n$.
\end{keyinsight}

The \texttt{Moments} trait abstracts over any data source that can produce
sample moments, so the same analytics work on slices, owned vectors, and
\texttt{PriceSeries}:

\begin{lstlisting}[language=Rust]
pub trait Moments {
    fn mean(&self) -> f64;
    fn variance(&self) -> Result<f64, CoreError>;
    fn std_dev(&self) -> Result<f64, CoreError>;
    fn skewness(&self) -> Result<f64, CoreError>;
    fn excess_kurtosis(&self) -> Result<f64, CoreError>;
}

impl Moments for &[f64] { /* delegate to free functions */ }
impl Moments for Vec<f64> { /* delegate to slice impl */ }
impl Moments for PriceSeries { /* delegate to slice impl */ }
\end{lstlisting}

\section{Rolling Windows}

\marginnote{%
\textbf{Rolling window}\\[0.3em]
For window $w$ on data of length $n$, the output has length $n - w + 1$ and
entry $i$ aggregates $\texttt{data}[i..i+w]$.
}

A rolling window applies an aggregation over each contiguous slice of length
$w$. The generic \texttt{rolling} function takes a closure, so any
aggregation (mean, std, max, custom) drops in:

\begin{lstlisting}[language=Rust]
pub fn rolling<F, T>(window: usize, data: &[f64], f: F)
    -> Result<Vec<T>, CoreError>
where F: Fn(&[f64]) -> T
{
    if window == 0 || window > data.len() {
        return Err(CoreError::InvalidWindow { window, len: data.len() });
    }
    let mut out = Vec::with_capacity(data.len() - window + 1);
    for i in 0..=(data.len() - window) {
        out.push(f(&data[i..i + window]));
    }
    Ok(out)
}

pub fn rolling_mean(window: usize, data: &[f64]) -> Result<Vec<f64>, CoreError> {
    rolling(window, data, mean)
}
\end{lstlisting}

\section{Hand-Rolled Simulation}

\subsection{xorshift64\*}

\marginnote{%
\textbf{xorshift64\*}\\[0.3em]
Three shifts and a multiply. Period $2^{64} - 1$. Fast, stateless, good
enough for Monte Carlo when cryptographic strength is not required.
}

We implement Vigna's xorshift64\* PRNG: three xor-shifts on a 64-bit state,
then a multiply by a constant. A seed of zero produces a degenerate stream,
so we replace it with a fixed default:

\begin{lstlisting}[language=Rust]
pub struct XorShift64 { state: u64 }

impl Rng for XorShift64 {
    fn next_u64(&mut self) -> u64 {
        let mut x = self.state;
        x ^= x >> 12;
        x ^= x << 25;
        x ^= x >> 27;
        self.state = x;
        x.wrapping_mul(0x2545_F491_4F6C_DD1D)
    }
    fn next_f64(&mut self) -> f64 {
        (self.next_u64() >> 11) as f64 / (1u64 << 53) as f64
    }
}
\end{lstlisting}

\subsection{Box-Muller Normal}

To sample $N(\mu, \sigma^2)$ without an external library, we use the
Box-Muller transform: two independent uniforms $u_1, u_2 \sim U(0,1)$ produce
one standard normal variate
\[
Z = \sqrt{-2 \ln u_1} \cos(2\pi u_2).
\]

\begin{lstlisting}[language=Rust]
pub fn box_muller_normal<R: Rng + ?Sized>(
    rng: &mut R, mean: f64, std: f64,
) -> f64 {
    let u1 = rng.next_f64().max(f64::EPSILON);
    let u2 = rng.next_f64();
    let r = (-2.0 * u1.ln()).sqrt();
    let theta = 2.0 * PI * u2;
    mean + std * (r * theta.cos())
}
\end{lstlisting}

The \texttt{Distribution} trait lets us plug in new distributions (exponential,
Student-$t$, jump-diffusion) without touching the simulator:

\begin{lstlisting}[language=Rust]
pub trait Distribution {
    fn sample<R: Rng + ?Sized>(&self, rng: &mut R) -> f64;
}

pub struct Normal { pub mean: f64, pub std: f64 }

impl Distribution for Normal {
    fn sample<R: Rng + ?Sized>(&self, rng: &mut R) -> f64 {
        box_muller_normal(rng, self.mean, self.std)
    }
}
\end{lstlisting}

\subsection{Geometric Brownian Motion}

\marginnote{%
\textbf{Exact GBM solution}\\[0.3em]
$S_{t+\Delta t} = S_t \exp\!\left(\left(\mu - \tfrac{1}{2}\sigma^2\right)\Delta t + \sigma\sqrt{\Delta t}\, Z\right)$\\[0.5em]
With $\sigma = 0$: $S_T = S_0 e^{\mu T}$, deterministic.
}

The stochastic differential equation
\[
dS_t = \mu S_t\, dt + \sigma S_t\, dW_t
\]
has the exact solution
\[
S_{t+\Delta t} = S_t \exp\!\left(\left(\mu - \tfrac{1}{2}\sigma^2\right)\Delta t + \sigma\sqrt{\Delta t}\, Z\right),
\]
where $Z \sim N(0, 1)$. We simulate one step at a time:

\begin{lstlisting}[language=Rust]
pub fn gbm_paths<R: Rng + ?Sized>(
    s0: f64, mu: f64, sigma: f64, t: f64,
    n_steps: usize, n_paths: usize, rng: &mut R,
) -> Vec<Vec<f64>> {
    let dt = t / n_steps as f64;
    let drift = (mu - 0.5 * sigma * sigma) * dt;
    let diffusion = sigma * dt.sqrt();
    let normal = Normal::standard();
    let mut paths = Vec::with_capacity(n_paths);
    for _ in 0..n_paths {
        let mut path = Vec::with_capacity(n_steps + 1);
        path.push(s0);
        let mut s = s0;
        for _ in 0..n_steps {
            let z = normal.sample(rng);
            s *= (drift + diffusion * z).exp();
            path.push(s);
        }
        paths.push(path);
    }
    paths
}
\end{lstlisting}

\begin{keyinsight}
With $\sigma = 0$, the $Z$ term drops out and $S_T = S_0 \exp(\mu T)$
deterministically. This is a useful sanity check: a zero-volatility GBM
should reproduce the risk-free growth exactly, with no RNG dependence.
\end{keyinsight}

\section{Fat Tails}

\begin{lstlisting}[language=bash]
$ cargo run -p quant-core --example fat_tails
1. Pure Gaussian (10k N(0,1) draws)
  skewness     = +0.0267
  excess kurt  = +0.0323

2. Fat tails (every 100th draw x10)
  skewness     = -0.3051
  excess kurt  = +74.1272

3. GBM terminal prices (10k paths, 252 steps)
  skewness     = +0.5957
  excess kurt  = +0.5959
\end{lstlisting}

The pure Gaussian sample has excess kurtosis near $0$, as expected. Scaling
every 100th draw by 10 injects rare large moves; the excess kurtosis jumps
from $\sim 0$ to $\sim 74$. GBM terminal prices are log-normally
distributed, so they have positive excess kurtosis (heavier right tail than
a Gaussian) and positive skewness --- a direct consequence of the
exponential in the exact solution.

\begin{keyinsight}
A Gaussian model fit to fat-tailed data will price a 4-sigma event as a
``1-in-15,000-years'' occurrence. The true frequency under fat tails may be
once a decade. This is why Value-at-Risk computed under Gaussian assumptions
systematically underestimates tail losses, and why option-implied volatilities
smile --- the market knows the tails are fatter than Gaussian allows.
\end{keyinsight}

\section{Exercises}

\begin{enumerate}
    \item \textbf{Sample vs population variance}: For
      $x \in \{1, 2, 3, 4, 5\}$, compute both the population variance
      (denominator $n$) and the sample variance (denominator $n - 1$). How
      large is the bias? When does it matter?

    \item \textbf{Skewness sign}: Construct a synthetic series with negative
      skew (long left tail) and one with positive skew. Verify
      \func{skewness} returns the correct sign. Which one is typical of
      equity returns?

    \item \textbf{RNG period}: The xorshift64\* generator has period
      $2^{64} - 1$. Write a test that verifies the first 1000 outputs from
      seed 42 are distinct, and reason about why ``distinct'' is necessary
      but not sufficient for a good PRNG.

    \item \textbf{GBM convergence}: Simulate 10\,000 GBM paths over one year
      (252 steps) with $S_0 = 100$, $\mu = 0.05$, $\sigma = 0.2$. Compute the
      mean terminal price. Compare with the analytical expectation
      $E[S_T] = S_0 e^{\mu T}$. How close is the Monte Carlo estimate?

    \item \textbf{Jump-diffusion}: Extend \func{gbm\_paths} with a Poisson
      jump term: with probability $\lambda \Delta t$ per step, multiply
      $S_t$ by $e^{J}$ where $J \sim N(\mu_J, \sigma_J^2)$. How does the
      excess kurtosis of terminal prices change as $\lambda$ grows?
\end{enumerate}